\section{Completion of the oracle-risk bound}
\label{sec:probability}

The final step uses a deleted graph only as a proof device. The
implemented procedure retains one full-graph eigensystem.
For each target $i$, put $B=A_{-i,-i}$ and recompute its degrees,
cap threshold, weights and projection $H_B$. Throughout this section, all constants
are uniform over the fixed structural class.

\begin{lemma}[Ordinary Bernoulli tilted cancellation]
\label{lem:sbm-floor-cancellation}
Let $w$ be a true block-constant Bernoulli log-odds vector on the
deleted graph and put
\[
 f=w-\bar w\one,\qquad \varepsilon=(H_B-I)f,
 \qquad \bar w=(n-1)^{-1}\one^\top w.
\]
On the deleted aggregate event, under the exact Chernoff tilt of
the raw-row oracle contrast, its Bernoulli probabilities $r_j$
satisfy
\begin{equation}\label{eq:sbm-floor-tilt}
 \left|\sum_{j\ne i}r_j\varepsilon_j\right|\le C,
 \qquad \sum_{j\ne i}r_j\varepsilon_j^2\le C.
\end{equation}
\end{lemma}

\begin{proof}
The tilted probabilities $r_j$ are constant within each true block
and satisfy $r_j\le Cp$. Write $\tau$ and $G$ for the threshold
and diagonal matrix recomputed from $B$. Since $r$ is block constant with $\norm r_\infty\le Cp$,
weighted orthogonality gives directly
\[
 r^\top\varepsilon
 =-\frac1\tau\big\langle(I-H_B)r,(I-H_B)f\big\rangle_G
       +\frac1\tau r^\top(\tau I-G)\varepsilon.
\]
The residual bounds give
$\norm{(I-H_B)r}_G\le Cp\sqrt n$ and
$\norm{(I-H_B)f}_G\le C\sqrt n$.
Also $\sum_j(G_j-\tau)^2/G_j\le Cn$, so
\[
 \norm{G^{-1/2}(\tau I-G)r}_2\le Cp\sqrt n,
 \qquad \norm\varepsilon_G\le C\sqrt n.
\]
Dividing either product by $\tau\asymp np$ gives $O(1)$.
This proves the first bound. Since $r_j\le Cp$ and
$\norm\varepsilon_2^2\le \frac{C}{p}$, the second bound follows as well.
The deterministic nonedge intercept remains in the oracle score
and is not absorbed into this cancellation.
When the cap is inactive on the deleted graph, $G=\tau I$ and
$H_B=\Pi_B=U_BU_B^\top$. The second term vanishes, leaving the
ordinary Euclidean adjacency identity
\[
 r^\top\varepsilon
 =-\big\langle(I-\Pi_B)r,(I-\Pi_B)f\big\rangle.
\]
\end{proof}

\begin{proof}[Proof of \cref{restore:sbm-theorem}]
Fix $A_0>0$. First consider ${np}\le A_0\log n$, with ${np}\to\infty$.
Let $\mathcal E$ collect the aggregate conditions
\eqref{restore:fit-geometry-one}--\eqref{restore:fit-block-event},
two-sided total degree, and their recomputed versions after deletion
of at most three vertices. For each prescribed fixed $r>0$,
\cref{prop:sparse-floor} gives these geometric conditions with
failure probability at most $C_rn^{-r}$, by increasing its
probability parameter to cover the polynomial number of minors.
For any true block total, independent-edge Bernstein gives
\[
 \PP\{|[Z^\top(A-{ZPZ^\top})Z]_{k\ell}|>Ln\}
 \le2\exp\{-cL^2n/({np}+L)\}.
\]
In the present range this is smaller than any fixed inverse power
of $n$, as is the total-degree failure. The same holds for the
required minors. For a target $i$, define $\mathcal E_i$ using
only its deleted graph and its further deletions. Then
\begin{equation}\label{eq:sbm-geometry-payment}
 \PP(\mathcal E^c)+\frac1n\sum_i\PP(\mathcal E_i^c)
 \le C_rn^{-r}.
\end{equation}

On $\mathcal E$, \cref{restore:fit-initializer,restore:fit-invariant}
give the initializer and every subsequent iterate the same bound
$\Mis(\widehat z,z)\le C/(np)$. At target
$i$ impose only $D_i\le L_D{np}$, its local spectral-row condition
\eqref{restore:fit-local-row}, and the requirement that the
incident row hits no exceptional deleted spectral row. By
\cref{restore:tail-energy-lemma,restore:tail-hit-corollary,restore:tail-target-corollary},
the average probability cost of these spectral exceptions is at
most $C_se^{-s{np}}$ for any prescribed fixed $s$, on choosing fixed
thresholds sufficiently large. Indeed, $G_j\ge\tau\asymp {np}$ gives
\eqref{restore:loo-hit} outside a deleted exceptional set
$\{j:\sqrt n\norm{(U_B)_{j\cdot}}_2>R\}$, whose expected hit
mass is controlled by \eqref{restore:tail-coefficient-count}.
The degree tail has the same bound by the single-row Bernoulli
Chernoff inequality. Polynomial factors in ${np}$ are absorbed by
first choosing a larger exponential rate. These are averages of
target probabilities, without a union over all tested vertices.

On the resulting local events, the deterministic fitted-score
bound and \cref{restore:loo-corrected} imply that an aligned error
at $i$ requires, for some $\ell\ne k=z_i$,
\begin{equation}\label{eq:sbm-independent-score}
 T_{k\ell}+\sum_{j\ne i}A_{ij}\varepsilon_j\le C,
 \qquad \varepsilon=(H_B-I)(w_{k\ell,-i}-\bar w_{k\ell,-i}\one).
\end{equation}
Here $T_{k\ell}$ is realized from this target's actual row, with its
full nonedge term and $n_k-1$ own-community trials. Its unconditional
law is the community-$k$ law defined in Section~\ref{sec:statements};
no independent substitute is introduced. The coefficients $\varepsilon$
are measurable from $B$ and independent of $(A_{ij})_{j\ne i}$. To apply the probability
bound, first use this deterministic inclusion, then drop the
full-graph event and condition only on $B$ on $\mathcal E_i$.
No independence is asserted conditional on $\mathcal E$.
The score threshold is fixed uniformly over all initialized refits,
because $1+np\Mis(\widehat z,z)$ is uniformly bounded by
\cref{restore:fit-invariant}.

The tilted mean and quadratic bounds in
\cref{lem:sbm-floor-cancellation} give the moment hypotheses of
\cref{lem:sparse-truncated-transfer}. Its exceptional-coefficient
hit term is at most $C_se^{-s{np}}$ by
\cref{restore:tail-hit-corollary}, after selecting a fixed
coefficient threshold. That truncation is used only in the proof.
Apply the sharp transfer to \eqref{eq:sbm-independent-score} using
the exact self-excluded row law. By the one-coordinate comparison in
Appendix~\ref{sec:sharp}, its exponential risk factor can be replaced
by $e^{-I_{k\ell}}$ at a uniformly bounded multiplicative cost.
Averaging targets and taking the fixed-$K$ union over competitors,
we obtain
\begin{equation}\label{eq:sbm-final-payment}
 \begin{split}
 \E\Mis(\widehat z,z)
 \lesssim{}&\sum_{k=1}^K\frac{n_k}{n}
                     \sum_{\ell\ne k}e^{-I_{k\ell}}\\
 &+e^{-snp}
       +\PP(\mathcal E^c)+\frac1n\sum_i\PP(\mathcal E_i^c).
 \end{split}
\end{equation}
Here the implicit factor may also depend on the chosen fixed
$s$ and the associated fixed local thresholds, never on $n$ or the
number of refits. We grouped the local testing terms by $k=z_i$;
the graph-dependent exceptional events retain their vertex indices.
The comparison in \eqref{eq:oracle-pairs} bounds the testing term by
a constant multiple of $R_{\mathrm{oracle}}$, for the specified tie rule.

On the compact structural class, choose a fixed $C$ such that
\[
 I_{k\ell}-\frac12\log(np)\le Cnp
\]
uniformly over the tested pairs. Consequently
\[
 R_{\mathrm{oracle}}\gtrsim\exp\{-Cnp-\tfrac12\log(np)\}.
\]
Choose $s>C+1$ and $r>A_0(C+1)+2$. When
$np\le A_0\log n$,
\[
 \sqrt{np}\left(e^{-(s-C)np}+n^{-r}e^{Cnp}\right)\longrightarrow0,
\]
which pays both remainders in \eqref{eq:sbm-final-payment}.

For ${np}>A_0\log n$, use \cref{restore:high-aggregate} in place of
the polynomially accurate geometry. It gives
$\PP(\mathcal E^c)+n^{-1}\sum_i\PP(\mathcal E_i^c)\le C_{s'}e^{-s'{np}}$ for any fixed
$s'>0$, with training events still measurable from the respective
deleted graphs. The row-tail proof uses deterministic aggregate
geometry and applies innovation concentration before restricting
to that event. Its estimates, the local LOO comparison and the
fitting argument therefore remain valid for this stronger event.
Equation \eqref{eq:sbm-final-payment} follows with the new geometry
remainder. Choosing $s,s'>C+1$ pays both exponential terms.
The two proof ranges cover every sequence ${np}\to\infty$ under
$np^2\le C_1$, proving the theorem without a logarithmic degree
lower bound.
\end{proof}
